\documentclass[a4paper]{article}

\newcommand{\docTitle}{Interaktive Systeme}

\usepackage{datetime2}  % for dates
\usepackage[utf8]{inputenc} % UTF-8 input encoding
\usepackage[T1]{fontenc} % fixes \textquotedbl with pdflatex/OT1
\usepackage{textcomp}   % additional text symbols for typewriter/upquote
\usepackage{xcolor}     % for gray color
\usepackage{listings}
\usepackage{marginnote} % for horizontal margin notes
\usepackage{parskip}
\usepackage{amsmath,amssymb,mathtools}
\usepackage{everypage}
\usepackage{fancyhdr}
\usepackage{amsfonts}   % Mengen and natural numbers
\usepackage{enumerate}  % Custom Enumerations
\usepackage[inline]{enumitem}   % List spacing and labels
\usepackage{multicol}   % Multiple Columns
\usepackage{tabularx}

\usepackage{hyperref}   % Links
\hypersetup{
    pdftitle={\docTitle},
    pdfauthor={Moritz Köhler},
    pdfkeywords={DHBW, Hochschule},
    colorlinks=true,
    linkcolor=black,
    filecolor=magenta,
    urlcolor=cyan
}

% -------
% Images
% -------

\usepackage{import}
\usepackage{xifthen}
\usepackage{pdfpages}
\usepackage{tikz}
\usetikzlibrary{arrows.meta, positioning}

\newcommand{\incfig}[1]{%
    \def\svgwidth{\columnwidth}
    \import{./fig/}{#1.pdf_tex}
}

% ---------------------------------
% Vertical Side note (Bottom Left)
% ---------------------------------

\newcommand{\verticaldocinfo}{%
  \rotatebox{90}{\tiny\color{gray}\texttt{\DTMtoday\_\docTitle\_\jobname}}%
}

\AddEverypageHook{
  \put(-40pt,-650pt){\verticaldocinfo}%
}

% ---------------
% Math Shortcuts
% ---------------

% Number sets
\newcommand{\R}{\mathbb{R}}
\newcommand{\N}{\mathbb{N}}
\newcommand{\Z}{\mathbb{Z}}
\newcommand{\Q}{\mathbb{Q}}
\newcommand{\C}{\mathbb{C}}

% Vectors and matrices
\newcommand{\vect}[1]{\mathbf{#1}}
\newcommand{\mat}[1]{\mathbf{#1}}

% Derivatives
\newcommand{\dd}{\mathrm{d}}
\newcommand{\dpartial}[2]{\frac{\partial #1}{\partial #2}}
\newcommand{\ddt}{\frac{\dd}{\dd t}}

% Norms and absolute values
\newcommand{\norm}[1]{\left\lVert #1 \right\rVert}
\newcommand{\abs}[1]{\left\lvert #1 \right\rvert}

% Probability & Statistics
\newcommand{\E}{\mathbb{E}}
\newcommand{\Var}{\mathrm{Var}}
\newcommand{\Prob}{\mathbb{P}}

% Fields and Algebras
\newcommand{\K}{\mathbb{K}}
\newcommand{\F}{\mathbb{F}}

% Set Operations
\newcommand{\compl}[1]{#1^{\mathrm{c}}}

% Matrix and Vector Operations
\newcommand{\T}{^{\top}}
\newcommand{\inv}{^{-1}}
\newcommand{\conj}[1]{\overline{#1}}

% -------------------
% Main lecture macro 
% -------------------

\newcommand{\lecture}[1]{%
  \protect\marginnote{\protect\tiny\protect\color{gray}#1}
}

% -----------------------------
% Command for box with heading
% -----------------------------

\fboxsep2mm
\newcommand{\know}[2]{
    \noindent\fbox{
        \parbox{\textwidth-21pt}{

            \textbf{#1:}
            \vspace{0.125cm}

            #2
        }
    }
}

% -------------------
% Listings Languages
% -------------------

% Shared defaults for code listings
\lstset{
    basicstyle=\ttfamily\small,
    keywordstyle=\color{blue}\bfseries,
    commentstyle=\color{gray},
    stringstyle=\color{teal},
    numberstyle=\tiny\color{gray},
    numbers=left,
    stepnumber=1,
    numbersep=8pt,
    showstringspaces=false,
    frame=single,
    breaklines=true,
    tabsize=4,
    keepspaces=false,
    columns=flexible,
    upquote=true,
    extendedchars=false % disable extended char handling to avoid UTF-8 conflicts
}

% SQL syntax highlighting
\lstdefinelanguage{SQL}{
    morekeywords={SELECT,FROM,WHERE,INSERT,INTO,VALUES,UPDATE,SET,DELETE,CREATE,TABLE,PRIMARY,KEY,NOT,NULL,AND,OR,JOIN,LEFT,RIGHT,INNER,ON,AS,ORDER,BY,GROUP,HAVING,LIMIT},
    sensitive=false,
    morecomment=[l]{--},
    morestring=[b]'
}

% JSON syntax highlighting
\lstdefinelanguage{json}{
    morestring=[b]",
    stringstyle=\color{teal},
    comment=[l]{//},
    morecomment=[s]{/*}{*/},
        morekeywords={true,false,null},
    literate=
     *{0}{{{\color{blue}0}}}{1}
      {1}{{{\color{blue}1}}}{1}
      {2}{{{\color{blue}2}}}{1}
      {3}{{{\color{blue}3}}}{1}
      {4}{{{\color{blue}4}}}{1}
      {5}{{{\color{blue}5}}}{1}
      {6}{{{\color{blue}6}}}{1}
      {7}{{{\color{blue}7}}}{1}
      {8}{{{\color{blue}8}}}{1}
      {9}{{{\color{blue}9}}}{1}
            {:}{{{\color{black}:}}}{1}
            {,}{{{\color{black},}}}{1},
    showstringspaces=false
}

% Language specific styles
\lstdefinestyle{javastyle}{
    language=Java,
    basicstyle=\ttfamily\small,
    keywordstyle=\color{blue}\bfseries,
    commentstyle=\color{gray},
    stringstyle=\color{teal},
    numberstyle=\tiny\color{gray},
    numbers=left,
    stepnumber=1,
    numbersep=8pt,
    showstringspaces=false,
    frame=single,
    breaklines=true,
    tabsize=4,
    keepspaces=true,
    columns=fullflexible,
    upquote=true
}

\lstdefinestyle{sqlstyle}{
    language=SQL,
    basicstyle=\ttfamily\small,
    keywordstyle=\color{blue}\bfseries,
    commentstyle=\color{gray},
    stringstyle=\color{teal},
    numberstyle=\tiny\color{gray},
    numbers=left,
    stepnumber=1,
    numbersep=8pt,
    showstringspaces=false,
    frame=single,
    breaklines=true,
    tabsize=4,
    keepspaces=true,
    columns=fullflexible,
    upquote=true
}

\lstdefinestyle{pythonstyle}{
    language=Python,
    basicstyle=\ttfamily\small,
    keywordstyle=\color{blue}\bfseries,
    commentstyle=\color{gray},
    stringstyle=\color{teal},
    numberstyle=\tiny\color{gray},
    numbers=left,
    stepnumber=1,
    numbersep=8pt,
    showstringspaces=false,
    frame=single,
    breaklines=true,
    tabsize=4,
    keepspaces=true,
    columns=fullflexible,
    upquote=true
}

\lstdefinestyle{cstyle}{
    language=C,
    basicstyle=\ttfamily\small,
    keywordstyle=\color{blue}\bfseries,
    commentstyle=\color{gray},
    stringstyle=\color{teal},
    numberstyle=\tiny\color{gray},
    numbers=left,
    stepnumber=1,
    numbersep=8pt,
    showstringspaces=false,
    frame=single,
    breaklines=true,
    tabsize=4,
    keepspaces=true,
    columns=fullflexible,
    upquote=true
}

\lstdefinestyle{bashstyle}{
    language=bash,
    basicstyle=\ttfamily\small,
    keywordstyle=\color{blue}\bfseries,
    commentstyle=\color{gray},
    stringstyle=\color{teal},
    numberstyle=\tiny\color{gray},
    numbers=left,
    stepnumber=1,
    numbersep=8pt,
    showstringspaces=false,
    frame=single,
    breaklines=true,
    tabsize=4,
    keepspaces=true,
    columns=fullflexible,
    upquote=true
}

\lstdefinestyle{jsonstyle}{
    language=json,
    basicstyle=\ttfamily\small,
    keywordstyle=\color{blue}\bfseries,
    commentstyle=\color{gray},
    stringstyle=\color{teal},
    numberstyle=\tiny\color{gray},
    numbers=left,
    stepnumber=1,
    numbersep=8pt,
    showstringspaces=false,
    frame=single,
    breaklines=true,
    tabsize=4,
    keepspaces=true,
    columns=fullflexible,
    upquote=true
}

% Default listing style
\lstset{language={}}

\title{\docTitle}
\author{Moritz}
\date{\today}

\begin{document}
\maketitle
\tableofcontents

\lecture{2026-10-01}

\section{Organisatorisch}

Wir machen meistens 2h Vorlesung und dann 1h Übungen

Seminararbeit wird am 26.11.26 veröffentlicht und am 08.01.27 abgegeben. AUFPASSEN das wir das Thema nicht mit KI verfehlen!

Fokus ist auf Interaktiv, nicht auf Systeme. Somit geht es primär um die Mensch Maschinen Interaktion.

\subsection{Überblick}

\begin{enumerate}
    \item Einführung (Motivation, Definition)
    \item Information (Menschliche Wahrnehmung, Gestaltwahrnehmung, Menschliche Informationsspeicherung, Hicks' Law und Fitts' Law)
    \item Interaktion ()
    \item Entwurfsprinzipien ()
    \item nicht mitgeschrieben
    \item nicht mitgeschrieben
    \item Usability Engineering ()
    \item (Hypertextsysteme ()) überspringen
    \item Barrierefreiheit ()
\end{enumerate}

\section{Einführung}

\subsection{Motivation}

Die Benutzerschnittstelle ist das einzige was ein Benutzer von einem System sieht. Ein System kann also nur so gut sein wie seine Schnittstelle.

Wir brauchen Design (First Buy Kriterium) und Usability (Second Buy Kriterium) als Marketingfaktor

Mittlerweile haben viele Gegenstände ein Interface, was bedient werden muss (Fernbedienung, Kaffeemaschine)

Jenach Nutzung gibt es andere Anforderungen. In der Fertigungshalle muss die Haltbarkeit und Austauschbarkeit leicht sein.

\subsubsection{Was möchte ein Benutzer von seinem "Gerät"?}

\begin{itemize}
    \item Das System darf keine Fehler machen.
    \item Wenn der Nutzer ein Fehler macht, soll das System dies korrigieren.
    \item Es soll nicht anstrengend sein.
    \item Bedienung soll Spaß machen
    \item viele weitere
\end{itemize}

Wenn auch nur ein Kriterium nicht erfüllt ist, wird sich das Produkt nicht durchsetzen!

\subsubsection{Mögliche Lösungsansätze}

\begin{itemize}
    \item User Centered Design
    \item Interaktionen mit dem zukünftigen Nutzer
    \item Lerne deine Benutzer kennen!
\end{itemize}

\subsubsection{Ein Benutzer}

\begin{itemize}
    \item Körperliches (Alter, Geschlecht, Gesundheit, Größe, ...)
    \item Soziales (Herkunft, Gruppenzugehörigkeit, Kulturelles Umfeld, Ausbildung, ...)
    \item Mentales (Erwartungen, Erfahrungen, Kognitive und Mentale Fähigkeiten, ...)
    \item Situatives (Tätigkeit, Arbeitsweise, Arbeitsabläufe, Benutzerhäufigkeit, Wichtigkeit der Aufgabe)
\end{itemize}

Welche Randbedingungen gibt es noch: Preis der Entwicklung, Verkaufspreis

% Hier fehlen Kleinigkeiten Mensch vs. Maschine und Auswirkungen von schlechtem Design

\subsubsection{Vorteile von einem "guten" Design}

Höhere Effizienz nach EN ISO 9241: Im Verhältnis zur Genauigkeit und Vollständigkeit eingesetzten Aufwand, mit dem Benutzer ein bestimmtes Ziel erreicht.

Höhere Produktivität

Weniger Benutzerfehler: Softwarefehler sind meist Design-Fehler. Gutes Design vermeidet und minimiert den Aspekt des "menschlichen Versagens"

Höhere Benutzerzufriedenheit evtl. nach EN ISO 9241.

Weniger Kosten: Nutzer braucht weniger Zeit (Mitarbeiter sind effiziente, Weniger Support, ...)

Wenn man die Regeln kennt, kann man schneller Gestalten.

Ein Vergleich:

Trägt jemand schlechte Kleidung, bemerkt man die Kleidung.

Trägt jemand gute Kleidung, achtet man auf die Person.

Dies lässt sich auch auf Systeme und ihre Schnittstelle übertragen. Beim Fehlen achtet man auf die Fehler, ohne Fehler achtete man auf das was dahinter steht.

\subsection{Definition}

User Interface, Benutzerschnittstelle (UI). Ist die technische Realisierung

Interaktives System (IS). Ist das UI was vom Nutzer benutzt und beeinflusst wird.

Human Computer Interaction (HCI). Definition nach ACM SIGCHI*. Wieder Interaktives System mit noch mehr wie Bewertung, Entwurf, Implementierung, usw.

Ergonomie: Wie meine Vertikale Maus.

Software Ergonomie: Anpassung von Software an die Menschlichen psychologischen Eigenschaften.

Benutzerfreundlichkeit, Gebrauchstauglichkeit (Usability). Erlebte Nutzungsqualität. Oder auch nach Norm DIN EN ISO 9241, Teil 11 - aufgeteilt nach Effektivität, Effizienz und Zufriedenheit.

\know{Seminararbeit}{Die Usability nach DIN EN ISO 9241, Teil 11 mit den drei Teilen ist wichtig.}

Universal Usability: Gilt für alle Nutzer. Sie müssen von allen Benutzer bedient werden können, wie ein Fahrkartenautomat.

Usability Engineering: Die Nutzerzentriete Entwicklung

User Experience (UX): Ist eine Erweiterung von Usability. \href{https://www.usability.de/usability-user-experience.html}{Vergleich Usability vs. UX}. Hier geht es um die gesamte Nutzungserlebnis vom Kauf, Nutzung bis zur Entsorgung/ Reparatur. Darauf achten dass alles was den Nutzer negativ beeinflusst zu vermeiden.

\end{document}